% Generated draft: requires mathematical and source review.
\documentclass[11pt]{article}
\usepackage[margin=25mm]{geometry}
\usepackage{fontspec}
\setmainfont{Latin Modern Roman}
\usepackage{amsmath,amssymb,mathtools,unicode-math}
\setmathfont{Latin Modern Math}
\usepackage{enumitem,needspace,xurl,hyperref}
\setlength{\parindent}{0pt}
\setlength{\parskip}{5pt}
\setlength{\emergencystretch}{3em}
\newcommand{\N}{\mathbb{N}}
\newcommand{\Z}{\mathbb{Z}}
\newcommand{\Q}{\mathbb{Q}}
\newcommand{\R}{\mathbb{R}}
\newcommand{\C}{\mathbb{C}}
\newcommand{\F}{\mathbb{F}}
\newcommand{\A}{\mathbb{A}}
\newcommand{\PP}{\mathbb{P}}
\newcommand{\E}{\mathbb{E}}
\newcommand{\eps}{\varepsilon}
\newcommand{\dd}{\,\mathrm{d}}
\begin{document}
\section*{UnsolvedMath ID: 20003330}
\textbf{Source problem code:} catalogue entry
\subsection{Definitions and mathematical statement}
The canonical record is the fragment ``Problem 40 for continuous spectrum was proved in [89,126]. Results toward''. It is a historical status sentence and unfinished transition, not the statement of Problem 40 or a new open problem. The aid identifies the context as continuous-spectrum quantum equidistribution: Luo--Sarnak's individual base-space compact-set ratio result for locally finite but globally infinite Eisenstein energy measures, and Jakobson's microlocal analogue using a positive Friedrichs-symmetrized distribution. The underlying space, spectral parameter and full theorem hypotheses are not supplied and are not reconstructed from the fragment.

The aid's independent normalization statement concerns a locally compact space $Z$, nonzero locally finite positive Radon measures $\mu_t$ and a nonzero target Radon measure $\nu$. Choose a nonnegative compactly supported continuous anchor $\phi_0$ with $\nu(\phi_0)>0$, and assume $\mu_t(\phi_0)>0$ eventually. Then the reported compact-test ratio convergence is
\[
 \frac{\mu_t(f)}{\mu_t(\phi_0)}
 \longrightarrow\frac{\nu(f)}{\nu(\phi_0)}
 \quad\text{for every }f\in C_c(Z).
\]
It is equivalent to vague convergence
\[
 \frac{\nu(\phi_0)}{\mu_t(\phi_0)}\,\mu_t\longrightarrow\nu.
\]
Every other compact nonnegative anchor with positive $\nu$ integral gives the same projective convergence notion. Positivity and nonzero anchor integrals are material. Multiplying an Eisenstein function by a nonzero scalar multiplies its energy measure by the scalar's squared modulus and cancels in these ratios. No normalization to a globally finite probability measure is implied when the energy measure has infinite total mass.

A cutoff estimate needs additional finite target data. Let $\nu$ now be a probability, let $0\leq\chi\leq1$ be compactly supported, and put $z=\nu(\chi)>0$. For bounded $f$, suppose the chosen scalar-normalized positive measure $\mu$ satisfies
\[
 |\mu(f\chi)-\nu(f\chi)|\leq\epsilon,\qquad
 |\mu(\chi)-z|\leq\delta<z.
\]
The reported error bound can be stated precisely as
\[
 \left|\frac{\mu(f\chi)}{\mu(\chi)}-\nu(f)\right|
 \leq\frac{\epsilon+\|f\|_\infty\delta}{z-\delta}
       +2\|f\|_\infty(1-z).
\]
The left side is a cutoff-normalized expectation versus the full target probability; it is not an expectation against the entire globally infinite source measure. Removing cutoffs requires control of numerator and denominator errors and of $1-z$.

The aid reports an explicit cusp half-line example in which fixed-cutoff limits recover the target volume while a diagonal growing cutoff loses a computable fraction of mass. Its actual measures, cutoff sequence and fraction are absent, so none is invented. The conclusion retained is that fixed-compact ratio convergence cannot be diagonalized without uniform control. Neither this normalization lemma nor the incomplete status sentence supplies a new proof of the full continuous-spectrum theorem or a formulation of the missing continuation.
\subsection{Short English statement}
This is historical context rather than a standalone question. The aid explains how one positive compact anchor normalizes projective vague convergence for possibly infinite energy measures. Its cutoff estimate needs explicit denominator control, and fixed-cutoff convergence alone does not justify a growing-cutoff limit.
\subsection{Sources}
\begin{itemize}
\item Local UnsolvedMath registry, permanent ID 20003330, source code catalogue entry; \texttt{aim-other-notes.json}, notes[136].
\item AIM, \emph{Emerging applications of measure rigidity}, status fragment about Problem 40: \url{https://aimath.org/WWN/measrigid/measrigid.pdf}. The fragment cites [89,126] but supplies neither Problem 40 nor full bibliography.
\item Luo--Sarnak and Jakobson continuous-spectrum equidistribution results are identified by the aid, the latter with a positive Friedrichs-symmetrized microlocal distribution. No complete theorem references are supplied.
\item Machine-generated research attempt 1: context only, not a problem at the run, rigor labelled ``proved,'' low novelty confidence; registry open/open-triage and reconstructed unverified statement. Not peer reviewed; independently verify before citation. Registry updated 14 August 2026.
\end{itemize}
\end{document}
